\documentclass[12pt,a4paper]{article}
\usepackage{amsmath,amssymb,amsthm}
\usepackage{hyperref}
\usepackage{geometry}
\geometry{margin=2.5cm}
\usepackage{enumitem}
\usepackage{booktabs}

\newtheorem{theorem}{Theorem}[section]
\newtheorem{lemma}[theorem]{Lemma}
\newtheorem{corollary}[theorem]{Corollary}
\newtheorem{proposition}[theorem]{Proposition}
\theoremstyle{definition}
\newtheorem{definition}[theorem]{Definition}
\theoremstyle{remark}
\newtheorem{remark}[theorem]{Remark}

\title{The Electroweak Scale from Recognition Science:\\
Hierarchy Dissolution via the $\varphi$-Ladder}

\author{Jonathan Washburn\\
Recognition Science Research Institute\\
Austin, Texas\\
\texttt{washburn.jonathan@gmail.com}}

\date{March 2026}

\begin{document}
\maketitle

\begin{abstract}
We derive the electroweak scale and dissolve the hierarchy
problem within Recognition Science (RS).  All mass scales in
RS are positions on the $\varphi$-ladder
($\varphi = (1+\sqrt{5})/2$), with no bare parameters.  The
Higgs vacuum expectation value
$v \approx 246~\mathrm{GeV}$ occupies rung
$r_v \approx 27$ relative to the electron mass, while the
Planck scale occupies rung $r_P \approx 107$.  The
electroweak--Planck hierarchy
$v/M_P \approx \varphi^{-80} \approx 10^{-17}$ is an
80-rung structural gap---a moderate integer determined by
the geometry of the $\varphi$-ladder, not a fine-tuning.
Radiative corrections do not destabilize the rung position
because RS has no bare parameters: masses are structural
outputs of the cost functional, not inputs subject to
renormalization.  No new particles (superpartners,
technifermions, KK excitations) are required.
\end{abstract}

%% ============================================================
\section{Introduction}
\label{sec:intro}
%% ============================================================

The hierarchy problem~\cite{tHooft1980,Weinberg1976} is the
question: why is the electroweak scale
$v \sim 10^2~\mathrm{GeV}$ so much smaller than the Planck
scale $M_P \sim 10^{19}~\mathrm{GeV}$?  In the Standard
Model (SM), the Higgs boson mass receives quadratically
divergent radiative corrections:
\begin{equation}
  \delta m_H^2 \sim \frac{\Lambda^2}{16\pi^2}
  \left(6y_t^2 - 3g^2 - g'^2 - 6\lambda\right),
  \label{eq:radiative}
\end{equation}
where $\Lambda$ is the UV cutoff.  If $\Lambda \sim M_P$,
the corrections are $\delta m_H^2 \sim
10^{38}~\mathrm{GeV}^2$, requiring a cancellation of 1 part
in $10^{34}$ against the bare Higgs mass-squared parameter
to achieve $m_H = 125~\mathrm{GeV}$~\cite{ATLAS2012}.

Proposed solutions include:
\begin{itemize}
  \item \textbf{Supersymmetry (SUSY)}~\cite{Dimopoulos1981}:
  Cancels quadratic divergences via superpartner loops.
  Requires $\geq 105$ new parameters.  No superpartners
  found at the LHC up to $\sim 2~\mathrm{TeV}$.
  \item \textbf{Composite Higgs}~\cite{Kaplan1984}: The Higgs
  is a pseudo-Nambu--Goldstone boson of a new strong sector.
  Constrained by precision electroweak data.
  \item \textbf{Extra dimensions}~\cite{ArkaniHamed1998}: The
  Planck scale is lowered by geometric dilution.
  Constrained by LHC searches.
  \item \textbf{Anthropic landscape}~\cite{Bousso2000}: The
  hierarchy is a selection effect among $\sim 10^{500}$
  string vacua.  Unfalsifiable.
\end{itemize}

RS~\cite{RS_Axioms} offers a different resolution: the hierarchy is
\emph{dissolved}, not solved.  There is no fine-tuning
because there are no bare parameters.

%% ============================================================
\section{Mass Scales on the $\varphi$-Ladder}
\label{sec:ladder}
%% ============================================================

\begin{definition}[$\varphi$-Ladder Positions]
\label{def:rungs}
In RS, every mass scale is a position on the
$\varphi$-ladder.  For any mass $m$, its rung position
relative to a reference mass $m_0$ is:
\begin{equation}
  r = \frac{\ln(m/m_0)}{\ln\varphi}.
\end{equation}
\end{definition}

\begin{proposition}[Electroweak and Planck Rungs]
\label{prop:rungs}
Taking the electron mass $m_e = 0.511~\mathrm{MeV}$ as
reference:
\begin{align}
  r_v &= \log_\varphi(v/m_e)
  = \frac{\ln(246{,}220/0.511)}{\ln(1.618)}
  \approx 27.2, \label{eq:rv} \\
  r_P &= \log_\varphi(M_P/m_e)
  = \frac{\ln(1.22 \times 10^{22}/0.511)}{\ln(1.618)}
  \approx 107.1. \label{eq:rP}
\end{align}
The hierarchy gap:
\begin{equation}
  \Delta r = r_P - r_v \approx 80~\text{rungs}.
  \label{eq:gap}
\end{equation}
\end{proposition}

\begin{theorem}[Hierarchy as Rung Gap]
\label{thm:hierarchy}
The electroweak--Planck hierarchy is:
\begin{equation}
  \boxed{\frac{v}{M_P} = \varphi^{-\Delta r}
  = \varphi^{-80} \approx 1.9 \times 10^{-17}.}
  \label{eq:hierarchy}
\end{equation}
\end{theorem}

\begin{remark}
The number $\Delta r = 80$ is a moderate integer, comparable
to other structural integers in the theory (e.g., $D = 3$,
$E = 12$, $V = 8$, $W = 17$).  There is nothing unnatural
about an 80-rung gap; it is the structural distance between
two positions on the ladder.
\end{remark}

%% ============================================================
\section{Dissolution of the Hierarchy Problem}
\label{sec:dissolution}
%% ============================================================

\begin{theorem}[No Bare Parameters]
\label{thm:no-bare}
In RS, the Higgs VEV is not a bare parameter plus radiative
corrections.  It is a structural output of the $J$-cost
minimization.
\end{theorem}

\begin{proof}
In the SM, the Higgs potential is:
\begin{equation}
  V(\phi) = -\mu^2 |\phi|^2 + \lambda |\phi|^4,
\end{equation}
where $\mu^2$ is the bare mass-squared parameter that
receives quadratic corrections from UV physics.  The
hierarchy problem is: why is $\mu^2_{\mathrm{phys}} \ll
\Lambda^2$?

In RS, the potential $V(\phi)$ \emph{emerges} from the
$J$-cost functional of the ledger.  The VEV
$v = \langle \phi \rangle$ is not an input parameter but
the solution to the cost minimization:
\begin{equation}
  \frac{\partial J_{\mathrm{EW}}}{\partial v} = 0
  \quad \Rightarrow \quad v \text{ at rung } r_v.
\end{equation}
Since $v$ is an output (the position where the $J$-cost
has its minimum), there is no ``bare value'' to protect
from radiative corrections.  The concept of fine-tuning is
inapplicable: one cannot fine-tune a structural output.
\end{proof}

\begin{remark}[Analogy with Atomic Energy Levels]
The hydrogen atom energy levels
$E_n = -13.6~\mathrm{eV}/n^2$ are structural solutions of
the Schr\"odinger equation.  They are not destabilized by
vacuum fluctuations (the Lamb shift provides a small
calculable correction, not a destabilization).  Similarly,
the $\varphi$-rung positions are structural solutions of the
cost minimization.
\end{remark}

\subsection{UV Finiteness of the Ledger}

\begin{proposition}[No Quadratic Divergences]
\label{prop:uv}
The recognition ledger is fundamentally discrete (8-tick
epochs, finite voxel structure).  There is no UV continuum
limit that generates quadratic divergences.
\end{proposition}

\begin{proof}
Quadratic divergences arise from integrating loop momenta
up to an arbitrarily high cutoff $\Lambda$:
\begin{equation}
  \int_0^\Lambda \frac{d^4k}{(2\pi)^4}\,
  \frac{1}{k^2 - m^2} \sim \Lambda^2.
\end{equation}
In RS, the ledger has a natural UV scale set by the tick
duration $\tau_0$ and the voxel spacing $\ell_0 = c\tau_0$.
Momenta above $k_{\max} = \pi/\ell_0$ do not exist in the
ledger.  The integral is naturally regulated:
\begin{equation}
  \int_0^{k_{\max}} \frac{d^4k}{(2\pi)^4}\,
  \frac{1}{k^2 - m^2} \sim k_{\max}^2 \sim \ell_0^{-2},
\end{equation}
which is finite and corresponds to the highest rung of the
$\varphi$-ladder.  There is no ``hierarchy'' between this
natural cutoff and the electroweak scale because the cutoff
is itself a rung position on the ladder.
\end{proof}

%% ============================================================
\section{Electroweak Masses from the VEV}
\label{sec:masses}
%% ============================================================

\begin{proposition}[Gauge Boson Masses]
\label{prop:gauge}
Given the VEV $v \approx 246~\mathrm{GeV}$, the
electroweak gauge boson masses follow from standard
electroweak symmetry breaking:
\begin{align}
  m_W &= \frac{g v}{2}
  \approx 80.4~\mathrm{GeV}, \\
  m_Z &= \frac{m_W}{\cos\theta_W}
  \approx 91.2~\mathrm{GeV},
\end{align}
where $g$ is the $SU(2)$ gauge coupling and $\theta_W$ is
the Weinberg angle (both derived from the $D = 3$ cube
geometry).
\end{proposition}

\begin{proposition}[Higgs Mass]
\label{prop:higgs}
The Higgs mass is:
\begin{equation}
  m_H = v\sqrt{2\lambda} \approx 125~\mathrm{GeV},
\end{equation}
where $\lambda$ is the quartic coupling determined by the
$J$-cost curvature at the VEV minimum.
\end{proposition}

\begin{remark}
The measured Higgs mass
$m_H = 125.25 \pm 0.17~\mathrm{GeV}$~\cite{ATLAS2012} is
consistent with the RS prediction.  The quartic coupling
$\lambda \approx 0.13$ is not a free parameter but is
determined by the cost functional.
\end{remark}

%% ============================================================
\section{Why No New Particles at the TeV Scale}
\label{sec:no-new}
%% ============================================================

\begin{theorem}[Desert Prediction]
\label{thm:desert}
RS predicts no new particles between the electroweak scale
and the Planck scale beyond those already in the Standard
Model.
\end{theorem}

\begin{proof}
In the SM with SUSY, new particles near the TeV scale
cancel quadratic divergences.  In RS, there are no
quadratic divergences (Proposition~\ref{prop:uv}), so no
cancellation mechanism is needed.  The $\varphi$-ladder
rungs between $r_v \approx 27$ and $r_P \approx 107$ are
occupied by the known SM particles (quarks, leptons, gauge
bosons) at their respective rungs---no new particle species
fill these intermediate rungs.
\end{proof}

\begin{remark}
The LHC null results for SUSY ($m_{\tilde{g}} >
2.3~\mathrm{TeV}$, $m_{\tilde{q}} >
1.8~\mathrm{TeV}$)~\cite{ATLAS_SUSY2023}, extra
dimensions, and composite Higgs particles are consistent
with the RS prediction.  Future colliders (FCC-hh at
$100~\mathrm{TeV}$) should continue to find no new
particles.
\end{remark}

%% ============================================================
\section{Comparison with Other Solutions}
\label{sec:comparison}
%% ============================================================

\begin{center}
\renewcommand{\arraystretch}{1.3}
\begin{tabular}{@{}lcccc@{}}
\toprule
\textbf{Approach} & \textbf{New} &
\textbf{Free} & \textbf{LHC} &
\textbf{$\Lambda$ role} \\
& \textbf{particles} & \textbf{params} &
\textbf{status} & \\
\midrule
SUSY & $\geq 12$ & $\geq 105$ & Not found &
Cancel \\
Composite Higgs & Yes & Several &
Constrained & Composite \\
Extra dims & Yes (KK) & Several &
Constrained & Dilute \\
Anthropic & No & $10^{500}$ vacua &
N/A & Select \\
RS $\varphi$-ladder & \textbf{No} & \textbf{0} &
\textbf{Consistent} & \textbf{Absent} \\
\bottomrule
\end{tabular}
\end{center}

\begin{remark}
The key distinction: SUSY, composite Higgs, and extra
dimensions \emph{solve} the hierarchy problem by
introducing new physics that cancels or screens the
quadratic divergence.  RS \emph{dissolves} it by removing
the conceptual framework (bare parameters + radiative
corrections) in which the problem is formulated.  The
anthropic approach also dissolves the problem, but is
unfalsifiable.  RS is falsifiable (see
\S\ref{sec:predictions}).
\end{remark}

%% ============================================================
\section{Predictions and Falsifiers}
\label{sec:predictions}
%% ============================================================

\begin{enumerate}
  \item \textbf{No SUSY particles}: No superpartners at any
  energy.  A confirmed discovery of a superpartner would
  indicate that the hierarchy is solved by SUSY rather than
  dissolved by the $\varphi$-ladder.

  \item \textbf{No composite Higgs signals}: The Higgs
  coupling deviations from SM predictions should remain
  below the percent level.  Large deviations would indicate
  compositeness.

  \item \textbf{No extra dimensions}: No KK excitations or
  microscopic black holes at any collider energy.

  \item \textbf{Standard Model desert}: Between the
  electroweak and Planck scales, no new particle species
  exist.  The ``great desert'' is a structural prediction.

  \item \textbf{Higgs mass stability}: The measured Higgs
  mass should remain consistent with the SM prediction
  from the $\varphi$-ladder, without additional corrections
  from new physics.

  \item \textbf{Vacuum stability}: The electroweak vacuum
  is stable (not metastable), because the $J$-cost minimum
  is unique and global.
\end{enumerate}

%% ============================================================
\section{Conclusion}
\label{sec:conclusion}
%% ============================================================

The hierarchy problem is dissolved in RS.  The electroweak
scale $v \approx 246~\mathrm{GeV}$ sits at rung
$r_v \approx 27$ on the $\varphi$-ladder; the Planck scale
at rung $r_P \approx 107$.  The 80-rung gap is a structural
feature of the ladder, not a fine-tuning.  No new particles
are required because there are no bare parameters to
protect---masses are structural outputs of the cost
minimization, not inputs subject to renormalization.  The
continuing null results at the LHC for SUSY, composite
Higgs, and extra dimensions are predictions of the
framework.

%% ============================================================
\begin{thebibliography}{99}

\bibitem{tHooft1980}
G.~'t Hooft, ``Naturalness, Chiral Symmetry, and
Spontaneous Chiral Symmetry Breaking,'' NATO Sci.\ Ser.\
B \textbf{59}, 135 (1980).

\bibitem{Weinberg1976}
S.~Weinberg, ``Implications of Dynamical Symmetry
Breaking,'' Phys.\ Rev.\ D \textbf{13}, 974 (1976).

\bibitem{ATLAS2012}
ATLAS Collaboration, ``Observation of a New Particle in
the Search for the Standard Model Higgs Boson,'' Phys.\
Lett.\ B \textbf{716}, 1 (2012).

\bibitem{Dimopoulos1981}
S.~Dimopoulos and H.~Georgi, ``Softly Broken
Supersymmetry and $SU(5)$,'' Nucl.\ Phys.\ B \textbf{193},
150 (1981).

\bibitem{Kaplan1984}
D.~B.~Kaplan and H.~Georgi, ``$SU(2) \times U(1)$ Breaking
by Vacuum Misalignment,'' Phys.\ Lett.\ B \textbf{136},
183 (1984).

\bibitem{ArkaniHamed1998}
N.~Arkani-Hamed, S.~Dimopoulos, and G.~Dvali, ``The
Hierarchy Problem and New Dimensions at a Millimeter,''
Phys.\ Lett.\ B \textbf{429}, 263 (1998).

\bibitem{RS_Axioms}
J.~Washburn, ``Recognition Science: Foundations,''
RS Axioms Document (2026).

\bibitem{Bousso2000}
R.~Bousso and J.~Polchinski, ``Quantization of Four-Form
Fluxes and Dynamical Neutralization of the Cosmological
Constant,'' JHEP \textbf{06}, 006 (2000).

\bibitem{ATLAS_SUSY2023}
ATLAS Collaboration, ``Search for Supersymmetry in Final
States with Missing Transverse Momentum and Multiple
$b$-jets,'' JHEP \textbf{2023}, 062 (2023).

\end{thebibliography}

\end{document}
